% !TEX root = ../main.tex

\section{Related Work}
\label{sec:related}

\noindent\textbf{Continual learning and distribution shift.}
A long line of work in machine learning addresses the failure of monolithic models under distribution drift. \emph{Continual learning} retains prior knowledge across changing tasks~\cite{}; \emph{domain adaptation} aligns source and target distributions~\cite{}; and \emph{test-time adaptation} updates the deployed model online from each batch~\cite{}. Mixture-of-experts methods take a complementary route, dispatching each input to a specialised sub-model rather than asking one model to fit all~\cite{}. All four operate on model weights. We adapt these intuitions to LLM agents at the \emph{harness} level (prompts, skills, tools, infrastructure), where weight updates are unavailable but harness mutations are; this shift gives the orthogonal regret decomposition $\Levo + \Lnav$ of \S\ref{sec:method}.

\noindent\textbf{Self-evolving agents.}
The closest LLM-agent precedents to our setting are self-evolving systems, which update the harness directly from execution feedback. A-Evolve~\cite{} introduces the basic linear-chain evolver: each cycle reads batch trajectories and mutates prompts, skills, memory, and tools. GEPA~\cite{} adds reflective Pareto prompt evolution using textual feedback rather than scalar reward. Meta-Harness~\cite{} uses a growing filesystem archive and Claude Code as the proposer. OctoTools~\cite{} is a training-free generalist Planner+Executor+Tool-Cards system intended as a hand-designed reference. All of these are primarily evaluated on i.i.d.\ benchmarks (HotpotQA, MATH, GAIA, MMLU-Pro, MedQA) where D1, D2, and D3 are absent or bounded, and all produce a single unconditional harness $\varphi(\Hist)$ rather than a task-conditional $\varphi(\Hist, x_t)$. We characterise the three task-intrinsic failure modes these systems exhibit on deployment-like streams, formalise the underlying mechanisms via the $\Levo + \Lnav$ decomposition, and introduce three architectural remedies that address each in turn.
